\section{Actionable Design for PTCTCL}
\subsection*{Step A: Instrument first}
Per level, emit:
\begin{itemize}
  \item touched chain count,
  \item pending-size distribution (median/p95/max),
  \item live-size distribution,
  \item top-k chain settle times,
  \item idle tail fraction near barrier.
\end{itemize}

\subsection*{Step B: Predictive ordering}
Build vector $V_L=\{(dk,\hat C)\}$ with
\[
\hat C = a\,|\Pend_{dk}|\,|\Live_{dk}| + b\,|\Pend_{dk}|^2,
\]
sort descending, schedule dynamically.

\subsection*{Step C: Heavy-chain chunking}
For chains with $|\Pend_{dk}| > \tau$, settle only $K$ pending entries, then requeue remainder.
This improves fairness and reduces tail without changing logical checks.

\subsection*{Step D: Adaptive policy}
At each level choose policy:
\begin{itemize}
  \item low skew: incremental path,
  \item high skew/high contention: fused path + predictive ordering + chunking.
\end{itemize}

\subsection*{Correctness note}
Any optimization must preserve the chain-local semantic constraints (subsumption and union behavior under accepted ordering rules).
Hence implementation should change scheduling granularity, not logical predicates.
